\chapter{Introduction}
\label{ch:introduction}

\section{Motivation}
Financial institutions are threatened less by ordinary daily fluctuations than by rare, large losses. Their
risk measures therefore concentrate on the tail of the loss distribution: Value at Risk (VaR), the loss level
exceeded with a given small probability, and Expected Shortfall (ES), the average loss beyond that level. Since
the Fundamental Review of the Trading Book, banks compute market-risk capital from the 97.5\% ES
\citep{bcbs2019}, and ES has the theoretical advantage of being a coherent risk measure
\citep{artzner1999,acerbi2002}.

Estimating these quantities is difficult because returns are not Gaussian \citep{mandelbrot1963,cont2001}.
Extreme moves occur far more often than a normal distribution predicts, volatility clusters in time, and
crashes such as March 2020 arrive as sudden jumps rather than as a sequence of ordinary moves. A model can fit
the centre of the return distribution well and still misjudge the worst one percent of outcomes, which is
exactly the part that VaR and ES measure.

Jump-diffusion models \citep{merton1976,kou2002} capture fat tails and crashes by adding a compound Poisson
jump to the diffusion. The density of returns under such a model solves a forward partial
integro-differential equation (PIDE), and both risk measures follow from that density: VaR is one of its
quantiles and ES an integral over its tail. The practical question is how to obtain the density accurately in
the tail, for a portfolio of several assets, and quickly enough to be used when market conditions change.

Classical methods each fail one of these requirements. Finite-difference solvers are accurate in one or two
dimensions but their cost grows exponentially with the number of assets. Monte Carlo scales to any dimension but
its tail estimates converge slowly, as $N^{-1/2}$ in the number of simulated paths, and must be repeated for every
new parameter set. Physics-informed neural networks (PINNs) \citep{raissi2019} offer a third route: a neural
network trained to satisfy the equation itself, which is mesh-free, can take model parameters as inputs, and,
once trained, is evaluated by a forward pass.

\section{Problem statement}
This thesis asks \emph{when a physics-informed neural network gives better tail-risk estimates than
conventional methods, and when its overhead outweighs the gain}. The aim is not to show that PINNs are better,
but to establish whether, when and why they are, with a design in which a negative answer is as informative as a
positive one. That requires an exact ground truth, so that every error is measured rather than estimated; strong
baselines, both numerical (finite differences, Monte Carlo) and statistical (GARCH, filtered historical
simulation, extreme value theory); and an honest accounting of cost.

\section{Research questions}
\begin{description}
  \item[RQ1 (density versus tail).] How accurately does the PINN recover the density in the bulk, and separately
    in the tail? A small average error can hide a large tail error.
  \item[RQ2 (risk accuracy and calibration).] How accurate are the resulting VaR and ES against exact values,
    and how well calibrated are they on real multi-asset data?
  \item[RQ3 (does the physics help?).] How does the PINN compare with the same network trained on data alone,
    and how many samples does the data-only network need to match it?
  \item[RQ4 (scaling with dimension).] How does accuracy and cost change with the number of assets, compared with
    finite differences (usable up to about three dimensions) and Monte Carlo?
  \item[RQ5 (latency).] Can one network, trained across a range of market parameters, answer risk queries faster
    than Monte Carlo at equal tail error, and after how many queries does its training cost pay back?
\end{description}

\section{Approach}
The study proceeds in stages of increasing difficulty, each validated before the next.
\begin{enumerate}
  \item A one-dimensional Merton model, where the exact density, a finite-difference solver and Monte Carlo
    agree to high precision. Here the PINN design is fixed and each of its components is tested by ablation.
  \item A $d$-asset model with one common market-jump factor, which keeps the jump integral one-dimensional and
    the exact density available in every dimension. The PINN is scaled from one to ten assets.
  \item An amortised PINN that takes the market parameters as inputs, so one trained network covers a whole
    range of market states, with a latency comparison against Monte Carlo.
  \item Daily one-day-ahead backtests on two five-asset baskets (NIFTY and S\&P sector ETFs, 2018--2026) against
    seven standard risk models, for a linear and an option portfolio.
  \item An ablation of the physics against a data-only network, and a set of controlled failure cases.
\end{enumerate}

\section{Contributions}
\begin{itemize}
  \item A PINN design for the forward jump-diffusion PIDE that learns the tail as accurately as the bulk: a
    log-density output, a log-form residual that measures relative error, Gauss--Hermite quadrature of the jump
    integral, and a normalisation penalty. Ablations quantify the contribution of each.
  \item A mesh-free extension to $d$ assets in which collocation, normalisation and risk all use one importance
    proposal built from the equation's coefficients, with an exact or stochastic (Hutchinson) trace. It reaches
    about 1\% ES error up to five assets and identifies ten assets as the breakdown point.
  \item An amortised network over volatility level, jump intensity and jump mean, with risk contributions and
    scenario analysis by differentiation, and a latency comparison that reports where Monte Carlo is faster.
  \item A real-data evaluation with Kupiec, Christoffersen, FZ0 and Diebold--Mariano tests in which the PINN's
    error is separated from the model's error by evaluating the same model with its exact density.
  \item A physics-versus-data ablation that measures the value of the equation in exact samples, a set of failure
    cases including a far-tail extrapolation failure and its remedy, and an open, tested and reproducible code
    base in which every figure and number is produced by a configuration-driven script.
\end{itemize}

\section{Scope}
The model is a jump-diffusion with constant coefficients, Gaussian jump sizes and one common jump factor;
time-varying volatility enters through a daily volatility scale. Results are reported for up to five assets,
with ten assets as the documented breakdown point. Stochastic volatility, idiosyncratic jumps and self-exciting
jump arrivals are left to future work.

\section{Outline}
\Cref{ch:literature} reviews tail-risk estimation, jump-diffusion models, physics-informed and other neural PDE
solvers, and identifies the gap. \Cref{ch:methodology} describes the model, the reference methods, the PINN
design and the evaluation protocol. \Cref{ch:results} reports the results stage by stage and answers the research
questions. \Cref{ch:conclusion} summarises the findings, their limitations and directions for further work.
